% Generated draft: requires mathematical and source review.
\documentclass[11pt]{article}
\usepackage[margin=25mm]{geometry}
\usepackage{fontspec}
\setmainfont{Latin Modern Roman}
\usepackage{amsmath,amssymb,mathtools,unicode-math}
\setmathfont{Latin Modern Math}
\usepackage{enumitem,needspace,xurl,hyperref}
\setlength{\parindent}{0pt}
\setlength{\parskip}{5pt}
\setlength{\emergencystretch}{3em}
\newcommand{\N}{\mathbb{N}}
\newcommand{\Z}{\mathbb{Z}}
\newcommand{\Q}{\mathbb{Q}}
\newcommand{\R}{\mathbb{R}}
\newcommand{\C}{\mathbb{C}}
\newcommand{\F}{\mathbb{F}}
\newcommand{\A}{\mathbb{A}}
\newcommand{\PP}{\mathbb{P}}
\newcommand{\E}{\mathbb{E}}
\newcommand{\eps}{\varepsilon}
\newcommand{\dd}{\,\mathrm{d}}
\begin{document}
\section*{UnsolvedMath ID: 6700079}
\textbf{Source problem code:} catalogue entry
\subsection{Definitions and mathematical statement}
A continuous Riemannian metric is a continuous positive-definite symmetric tensor; a $C^2$ metric has classical scalar curvature. The source uses a volumic scalar-curvature quantity $\operatorname{Sc}^{\mathrm{vol}}_n(g)$ for continuous metrics without supplying its definition.

The proposed topological-equivalence conjecture is: if a smooth $n$-manifold $M$ admits a continuous metric $g_1$ with $\operatorname{Sc}^{\mathrm{vol}}_n(g_1)>0$, then it admits a $C^2$ metric $g_2$ with
\[
\operatorname{Sc}_{g_2}(x)>0\qquad(x\in M).
\]
The extract prints $\operatorname{Sc}(g_1)>0$ in the conclusion after introducing $g_2$; reading the conclusion as scalar positivity of $g_2$ is an explicitly flagged correction, not a verified transcription. No approximation or closeness of the two metrics is requested.

For closed $M$, the source offers a second formulation: if there is no smooth positive-scalar-curvature metric, then every continuous metric $g$ has balls at arbitrarily small radii whose volume is at least Euclidean. Writing $\omega_n=\operatorname{Vol}(B^n_{\mathrm{Eucl}}(1))$, its natural quantified reading is
\[
\forall\varepsilon>0\ \exists x\in M\ \exists R\in(0,\varepsilon):
\operatorname{Vol}_g(B_g(x,R))\geq\omega_n R^n.
\]
The centers may vary with the radius. The source omits the volume symbol on the left of its comparison; the displayed formula interprets it as ball volume. Without the volumic-curvature definition, the asserted equivalence between this small-ball condition and the first conjecture cannot be verified. The first sentence is not explicitly restricted to closed manifolds, whereas this restatement is.
\subsection{Short English statement}
Does existence of a continuous metric with positive volumic scalar curvature imply existence of a $C^2$ metric with positive classical scalar curvature? The volumic definition is missing, and the extract has a metric-label typo and an incompletely justified small-ball reformulation.
\subsection{Sources}
\begin{itemize}
\item Misha Gromov, \emph{101 Questions, Problems and Conjectures around Scalar Curvature} (2017), printed tag [?81], occurrence 82, PDF page 82, ``Topological Equivalence of Different Scalar Curvatures,'' \url{https://www.ihes.fr/~gromov/wp-content/uploads/2018/08/101-problemsOct1-2017.pdf}; accessed 2026-07-29, PDF-text-normalized extraction. The provenance describes the public author PDF as incomplete.
\item The final sentence refers to a probably unrealistic strengthening of the $\mathbb Q$-non-essentiality conjecture, but its statement and reference are truncated at [?; no missing strengthening is reconstructed.
\item Local registry \texttt{data/unsolvedmath/problems.json}, permanent UnsolvedMath ID $6700079$, source code catalogue entry, \emph{Scalar Curvature Question [?81]}. Dataset attribution: Ulam AI and the UnsolvedMath Contributors.
\end{itemize}
\end{document}
